%==============================================================================
%  ملحق - الذكاء الاصطناعي والتحليل التنبّئي
%==============================================================================

%==============================================================================
\chapter{ملحق --- الذكاء الاصطناعي والتحليل التنبّئي}
\label{chap:ia}
%==============================================================================
يجيب هذا الملحق عن أسئلة \emph{لماذا} و\emph{كيف} و\emph{أين} الذكاء الاصطناعي
في \textbf{Zümm}. لا يُنشئ أي متطلب جديد: يُمدّد شقّ المستشعرات والتحليل التنبّئي
المُستبَقَين في الفصل~\ref{chap:conception} (§\,4.5) ومسيرة الخدمة المصغّرة
\emph{للمعطيات} المذكورة في الملحق~\ref{chap:migration}. الهدف توثيق دمج للذكاء
الاصطناعي \textbf{متّسق مع البنية ضعيفة الاقتران}، \textbf{دون خيانة قيود} كراس
الشروط (النشر الذاتي، المجانية، التحقّق البشري).

\begin{encadre}
\textbf{التموضع:} في Zümm، الذكاء الاصطناعي \textbf{وحدة مساعدة على القرار}، لا
عضو تحكّم أبدًا. يُنتج \emph{تنبيهات مسبقة} قابلة للتفسير؛ يحتفظ عون تربية النحل
بـ\textbf{المصادقة النهائية} عبر تقرير الزيارة، وفق قاعدة التقييم في §\,4.5.4.
\end{encadre}

%------------------------------------------------------------------------------
\section{المبادئ الموجِّهة}
%------------------------------------------------------------------------------
تحترم كل لبنة ذكاء اصطناعي معتمَدة أربعة مبادئ مستمدّة مباشرةً من متطلبات المشروع:
\begin{itemize}
  \item \textbf{المساعدة على القرار والتحقّق البشري}: الذكاء الاصطناعي يُشير، لا يقرّر. كل مُخرَج يؤكّده العون (§\,4.5.4).
  \item \textbf{النشر الذاتي دون خدمة خارجية إجبارية}: لا تعتمد أي وظيفة أساسية على واجهة سحابية مدفوعة. تعمل المعالجات الأساسية \textbf{محليًا}؛ ويبقى اللجوء إلى السحابة \emph{اختياريًا} وقابلًا للتفعيل بالإعداد.
  \item \textbf{الفصل والنمطية}: الذكاء الاصطناعي \textbf{وحدة قابلة للتفعيل} عبر \texttt{ConfigZumm.ini}، معزولة عن نواة الأعمال بواجهة، بروح الاقتران الضعيف (الملحق~\ref{chap:solid}).
  \item \textbf{القابلية للتفسير والاقتصاد}: تُفضّل الطرق \textbf{القابلة للتفسير} (إحصاء، قواعد) والمُقتصِدة، بدل النماذج المُعتِمة، للبقاء قابلًا للدفاع وللاستنساخ.
\end{itemize}

%------------------------------------------------------------------------------
\section{خريطة استعمالات الذكاء الاصطناعي}
%------------------------------------------------------------------------------
يُحصي الجدول أدناه استعمالات الذكاء الاصطناعي المناسبة لـ Zümm، مُميّزًا ما هو
\textbf{ضمن النطاق} (قابل للتنفيذ في الأساس الحالي) عمّا هو \textbf{مُستبَق}
(واجهة مُحدَّدة، تطوير لاحق ضمن شقّ المستشعرات).

\begin{xltabular}{\textwidth}{>{\bfseries}p{3.1cm} X C{1.7cm} C{1.6cm} C{1.9cm}}
\toprule
\textbf{الاستعمال} & \textbf{الوصف} & \textbf{تبعية سحابية} & \textbf{الجهد} & \textbf{الحالة} \\
\midrule
\endhead
كشف شذوذ تكيّفي &
خطّ أساس لكل خلية يحلّ محلّ العتبة الثابتة، لكشف انحراف غير معتاد في الإنتاج أو
العدد أو الحرارة. &
لا & منخفض & \textbf{مُعتمَد} \\
كشف صوتي للتطريد &
تحليل البصمة الصوتية لتنبيه مسبق بالتطريد من 1 إلى 5 أيام مقدّمًا (تربية النحل
الدقيقة، §\,4.5.1). &
لا & مرتفع & مُستبَق \\
رؤية: كشف الأمراض &
مساعدة على التحديد (الفاروا، الحضنة) انطلاقًا من صور التفتيش المرفقة بتقرير
الزيارة. &
لا & مرتفع & مُستبَق \\
إدخال صوتي مُساعَد &
إملاء المحضر، والنسخ يغذّي حقول التقرير (§\,4.7، أولوية منخفضة). &
اختياري & متوسّط & مُستبَق \\
مساعد التلخيص &
تلخيص بلغة طبيعية للوحات القيادة (العائد، التنبيهات) للمالك. &
اختياري & متوسّط & اختياري \\
\bottomrule
\end{xltabular}

\begin{encadre}
\textbf{توصية بالأولوية:} تسليم \textbf{كشف الشذوذ التكيّفي} كوظيفة حقيقية (قابل
للتنفيذ في الأساس المفروض، دون تبعية سحابية)، و\textbf{تحديد الاستعمالات الأخرى
كواجهات مُستبَقة}. هذا هو الموقف الذي اعتمده كراس الشروط سلفًا لشقّ المستشعرات،
ما يحفظ اتّساق الوثيقة.
\end{encadre}

%------------------------------------------------------------------------------
\section{الحالة المُعتمَدة: كشف الشذوذ التكيّفي}
%------------------------------------------------------------------------------
يعرّف §\,4.5.4 \textbf{عتبات ثابتة} (مثلًا الحرارة الداخلية أقل من
32\,\textdegree م أو أكثر من 36\,\textdegree م) وقاعدة مانع ارتداد. حدّ العتبة
الثابتة أنها تتجاهل \textbf{السلوك الخاصّ لكل خلية} و\textbf{موسميتها}. التطوّر
الطبيعي هو تعلّم \textbf{خطّ أساس لكل خلية} وقياس الانحراف عنه، ما يكشف «الانخفاض
النسبي أكثر من 20\,\% مقارنةً بالفترة السابقة» بشكل تكيّفي.

\subsection{الصياغة (دون شيفرة)}
لكل زوج (خلية، مؤشّر)، يُحتفَظ بـ\textbf{متوسّط متحرّك أُسّي} (EWMA) $\mu_t$
وتقدير تشتّت $\sigma_t$، يُحدَّثان عند كل قياس مُسوًّى جديد $x_t$ (الوحدة
المرجعية، §\,5.3):
\begin{align}
  \mu_t      &= \alpha\,x_t + (1-\alpha)\,\mu_{t-1}, \\
  v_t        &= \alpha\,(x_t-\mu_{t-1})^2 + (1-\alpha)\,v_{t-1},
                \qquad \sigma_t = \sqrt{v_t}.
\end{align}
\textbf{درجة الشذوذ} هي الانحراف المُسوَّى (\emph{z-score} منزلق):
\begin{equation}
  z_t = \frac{\lvert x_t - \mu_{t-1}\rvert}{\sigma_{t-1} + \varepsilon}.
\end{equation}
لا يُطلَق \textbf{تنبيه مسبق} إلا إذا $z_t > k$ (الحساسية) \textbf{باستمرار على
$N$ قياسات متتالية} (تخلفية، كما في §\,4.5.4). هذه الاستمرارية مع القاعدة القائمة
مقصودة: تُستبدَل عتبة مطلقة بعتبة \emph{نسبية مُتعلَّمة}، دون تغيير منطق مانع
الارتداد ولا مبدأ التحقّق البشري.

\subsection{المعلمات (كلها قابلة للضبط عبر ConfigZumm.ini)}
\begin{xltabular}{\textwidth}{>{\bfseries}p{2.6cm} p{2.4cm} X}
\toprule
\textbf{المعلمة} & \textbf{الدور} & \textbf{أثر الضبط} \\
\midrule
\endhead
$\alpha$ (التنعيم) & وزن القياس الحديث & كلما كبر $\alpha$، أسرع تفاعل خطّ الأساس (وأكثر حساسيةً للضجيج). \\
$k$ (الحساسية) & عتبة على z-score & كلما صغر $k$، أكثر حساسية الكشف (وأكثر التنبيهات الكاذبة). \\
$N$ (الاستمرار) & عدد القياسات المتتالية & كلما كبر $N$، أقلّ التنبيهات الكاذبة (بثمن تأخّر طفيف). \\
النافذة / الموسم & التاريخ المرجعي & يتيح خطّ أساس لكل موسم (المِيلة، التشتية). \\
\bottomrule
\end{xltabular}

\begin{figure}[htbp]
\centering
\includegraphics[width=0.72\linewidth]{diagramme_activite_ia_pipeline}
\caption{خطّ أنبوب كشف الشذوذ التكيّفي، من القياس الخام إلى التنبيه المسبق
المؤكَّد من العون.}
\label{fig:ia-activite}
\end{figure}

\begin{encadre}
\textbf{الاتّساق مع §\,4.5.4:} تُحوَّل القيمة أولًا \textbf{إلى الوحدة المرجعية}
(Strategy التحويل، الملحق~\ref{chap:solid})، ثم تُقارَن بعتبة؛ يبقى التنبيه
خاضعًا لـ\textbf{مانع الارتداد} ويظلّ \textbf{مساعدةً على القرار}. تتغيّر طبيعة
العتبة فقط: من \emph{ثابتة} إلى \emph{مُتعلَّمة لكل خلية}.
\end{encadre}

%------------------------------------------------------------------------------
\section{بنية الدمج}
%------------------------------------------------------------------------------
يندمج الذكاء الاصطناعي دون كسر أيّ شيء في البنية الطبقية
(الشكل~\ref{fig:couches}). تُستعمَل آليتا فصل.

\subsection{الكاشف كـ Strategy قابلة للتبادل}
يصبح كشف الشذوذ \textbf{استراتيجية لـ \texttt{RegleAlerte}} (نمط Strategy،
الملحق~\ref{chap:solid})، شأنه شأن العتبة الثابتة القائمة. يختار
\texttt{GestionnaireAlertes}، بالإعداد، بين «عتبة ثابتة» و«خطّ أساس تكيّفي» دون
تغيير شيفرته (مبدأ الانفتاح/الانغلاق). فالذكاء الاصطناعي ليس إضافةً اقتحامية، بل
\textbf{تنفيذًا إضافيًا لواجهة مُقرَّرة سلفًا}.

\begin{figure}[htbp]
\centering
\includegraphics[width=\linewidth]{diagramme_sequence_ia_detection}
\caption{تسلسل وحدة الذكاء الاصطناعي: يُقيَّم القياس المُستوعَب بالاستراتيجية
التكيّفية، التي تُطلق تنبيهًا مسبقًا يؤكّده العون.}
\label{fig:ia-sequence}
\end{figure}

\subsection{المعالجات الثقيلة كخدمة مصغّرة مفصولة}
الاستعمالات التي تتطلّب \textbf{إشارة أو رؤية} (صوت، صور) يجب \textbf{ألا}
تُنفَّذ أبدًا في حاوية EJB. تُرحَّل إلى \textbf{خدمة مصغّرة بلغة Python} (pandas،
scikit-learn / TensorFlow)، مكشوفة للنواة Java عبر \textbf{REST/JSON}، وهي
بالضبط مسيرة \emph{المعطيات} الموصوفة في الملحق~\ref{chap:migration}. هذه الخدمة
المصغّرة \textbf{اختيارية}: في غيابها، يعمل الشقّ اليدوي والكشف الإحصائي الأساسي
بالكامل.

\begin{figure}[htbp]
\centering
\begin{tikzpicture}[
  noeud/.style={draw=mielfonce,fill=grisdoux,thick,minimum width=3.1cm,minimum height=1.1cm,font=\footnotesize,align=center},
  ianode/.style={draw=mielfonce,fill=miel!15,thick,minimum width=3.1cm,minimum height=1.1cm,font=\footnotesize,align=center},
  fl/.style={-{Latex[length=2mm]},mielfonce,thick}
]
\node[noeud] (capt) {بوّابة\\ميدانية};
\node[noeud,right=1.5cm of capt] (coeur) {نواة Zümm (J2EE)\\Servlet + EJB + JDBC};
\node[noeud,right=1.5cm of coeur] (bd) {PostgreSQL\\سلاسل زمنية};
\node[ianode,below=1.1cm of coeur] (ia) {خدمة مصغّرة للذكاء الاصطناعي (Python)\\pandas / scikit-learn};
\node[noeud,above=1.0cm of coeur] (agent) {عون\\لوحة قيادة};
\draw[fl] (capt)--(coeur) node[midway,above,font=\scriptsize]{REST / MQTT};
\draw[fl] (coeur)--(bd) node[midway,above,font=\scriptsize]{JDBC};
\draw[fl] (coeur.south)--(ia.north) node[midway,right,font=\scriptsize]{REST/JSON};
\draw[fl] (ia.north west) to[bend left=15] node[midway,left,font=\scriptsize]{score} (coeur.south west);
\draw[fl] (coeur)--(agent) node[midway,right,font=\scriptsize]{تنبيه مسبق};
\node[draw=mielfonce,dashed,thick,inner sep=6pt,fit=(coeur)(bd),label={[font=\scriptsize\bfseries,mielfonce]above:الأساس المفروض (J2EE)}] {};
\end{tikzpicture}
\caption{دمج مفصول: يبقى الأساس J2EE مُتحكّمًا، والخدمة المصغّرة للذكاء الاصطناعي
وحدة اختيارية مُوصَّلة عبر خدمة ويب.}
\label{fig:ia-archi}
\end{figure}

%------------------------------------------------------------------------------
\section{الاستعمالات المُستبَقة (واجهات مُحدَّدة، خارج نطاق التطوير)}
%------------------------------------------------------------------------------
وفق معالجة شقّ المستشعرات، تُحدَّد الاستعمالات المتقدّمة \textbf{كواجهات} كي
يستوعبها نموذج المعطيات والبنية، دون تطويرها في هذا المشروع.
\begin{itemize}
  \item \textbf{الصوت / التطريد}: يستوعب نموذج المعطيات سلفًا مؤشّر \texttt{SIGNATURE\_ACOUSTIQUE} (§\,4.5). تُرجع الخدمة المصغّرة حكمًا (نمط تطريد مُكتشَف، وجود الملكة) يغذّي تنبيهًا مسبقًا.
  \item \textbf{الرؤية / الأمراض}: تشكّل \textbf{صور التفتيش} المرفقة سلفًا بتقرير الزيارة مجموعة الدخل؛ يُرجع نموذج مُدرَّب مسبقًا اقتراحًا، يؤكّده العون.
  \item \textbf{الإدخال الصوتي}: يمكن أن يعتمد النسخ على \textbf{Web Speech API للمتصفّح} (مجانية، من جهة العميل، دون تبعية للخادم)، مع تراجع سحابي \emph{اختياري} مُفعَّل بـ \texttt{ConfigZumm.ini}.
\end{itemize}

%------------------------------------------------------------------------------
\section{الأخلاقيات والحدود والامتثال}
%------------------------------------------------------------------------------
\begin{xltabular}{\textwidth}{>{\bfseries}p{4.4cm} X}
\toprule
\textbf{القضية} & \textbf{الترتيب المعتمَد} \\
\midrule
\endhead
تنبيهات كاذبة & تخلفية على $N$ قياسات وحساسية $k$ قابلتان للضبط؛ يبقى التنبيه إرشاديًا. \\
الإفراط في الثقة بالأتمتة & تموضع صريح كمساعدة على القرار؛ يحتفظ العون بالمصادقة النهائية (§\,4.5.4). \\
القابلية للتفسير & طرق قابلة للتفسير (درجة + عتبة) مُفضَّلة على النماذج المُعتِمة؛ يعرض التنبيه المسبق تبريره. \\
حماية المعطيات & معالجة محلية افتراضيًا، تشفير SSL / X.509 للتبادلات، لا نقل سحابي مفروض. \\
غياب الاحتكار & الخدمة المصغّرة للذكاء الاصطناعي اختيارية وقابلة للاستبدال؛ بدونها، تبقى الإدارة والكشف الأساسي فاعلَين. \\
الاقتصاد & أولوية للمعالجات الإحصائية الخفيفة؛ لا يُطلَق الاستدلال الثقيل إلا عند الطلب. \\
\bottomrule
\end{xltabular}

\begin{encadre}
\textbf{الخلاصة:} يتلخّص ذكاء Zümm الاصطناعي في قاعدة واضحة --- \emph{تعلّم
الحالة الطبيعية لكل خلية، الإشارة إلى الانحراف، ترك القرار للإنسان}. تضيف هذه
المقاربة قيمة تنبّئية مع احترام كامل لقيود النشر الذاتي والاقتران الضعيف والتحقّق
البشري في كراس الشروط. تحوّل الشقّ التنبّئي \emph{المُستبَق} إلى مسيرة تنفيذ
\textbf{ملموسة وتدريجية}.
\end{encadre}
